\documentclass[11pt]{article}

\usepackage[margin=1in]{geometry}
\usepackage{amsmath}
\usepackage{amssymb}
\usepackage{booktabs}
\usepackage{array}
\usepackage{hyperref}

\title{Notes On The Viscosity Models Used In The ORCA Flow-Log Review}
\author{HFE System Analysis Notes}
\date{\today}

\begin{document}

\maketitle

\section{Purpose}

This note explains the viscosity models currently used in the ORCA flow-log
review workflow for the stable cooldown hold. The goal is not to claim a
universal constitutive law for the fluid. Instead, the goal is to describe the
practical local models used to compare measured cooldown data against a single
reference curve derived from published 3M viscosity data.

The models discussed here are:
\begin{enumerate}
    \item the local smooth law derived from 3M measurements,
    \item the density-power fit,
    \item the flow-only proxy, now based on volumetric flow,
    \item the combined pump-flow law,
    \item and the smoothed combined-law trendline used only for visualization.
\end{enumerate}

\section{Notation}

Throughout this document:
\begin{itemize}
    \item $T$ is the measured flow-meter temperature in degrees Celsius,
    \item $\rho$ is the measured liquid density in $\mathrm{kg/m^3}$,
    \item $Q$ is the measured volumetric flow in $\mathrm{L/min}$,
    \item $\nu$ is the kinematic viscosity in $\mathrm{cSt}$,
    \item $\mu$ is the dynamic viscosity in $\mathrm{cP}$.
\end{itemize}

The kinematic and dynamic viscosities are related by
\begin{equation}
    \mu = \nu \frac{\rho}{1000},
\end{equation}
because $1~\mathrm{cSt} \cdot 1~\mathrm{kg/m^3} / 1000 = 1~\mathrm{cP}$.

\section{Reference Model: Local Smooth Law From 3M Data}

The reference model is constructed from the published 3M viscosity table for
the working fluid. Rather than use the full table directly, a smooth local fit
is built over the temperature window that is most relevant to the measured
cooldown segment.

The fitted form is
\begin{equation}
    \ln \nu_{\mathrm{ref}}(T) = a_2 T^2 + a_1 T + a_0,
\end{equation}
so that
\begin{equation}
    \nu_{\mathrm{ref}}(T) = \exp\!\left(a_2 T^2 + a_1 T + a_0\right).
\end{equation}

This model is called a \emph{local smooth law} because:
\begin{itemize}
    \item it is anchored to the published 3M values,
    \item it is smooth rather than piecewise,
    \item and it is only intended for the measured temperature window, not for
          broad extrapolation.
\end{itemize}

This is the primary reference curve. The other models are judged by how well
they reproduce this local reference during the stable cooldown hold.

\section{Density-Power Fit}

\subsection{Definition}

The density-power fit assumes that, over the narrow measured range,
kinematic viscosity can be approximated by a power law of density:
\begin{equation}
    \ln \nu_{\rho} = A + B \ln \rho,
\end{equation}
or equivalently
\begin{equation}
    \nu_{\rho} = C \rho^B,
    \qquad
    C = e^A.
\end{equation}

The coefficients $A$ and $B$ are not fitted to a direct viscosity measurement.
Instead, they are fitted so that $\nu_{\rho}$ matches the local 3M-derived
reference viscosity $\nu_{\mathrm{ref}}(T)$ at the temperatures and densities
observed in the stable cooldown hold.

\subsection{Interpretation}

This fit is best understood as a \emph{local empirical surrogate}. It works
well when:
\begin{itemize}
    \item the fluid stays in a single liquid phase,
    \item pressure is roughly constant,
    \item and $\rho(T)$ is nearly one-to-one over the measured interval.
\end{itemize}

Under those conditions, density acts as a strong proxy for thermal state, so
viscosity can be reconstructed from density alone with good local accuracy.

\subsection{Scientific Status}

The density-power fit should not be presented as a universal fluid law. There
is no general theorem that viscosity must follow a single power of density over
all temperatures and pressures. The model is scientifically acceptable as a
local interpolation tool, but not as a broad constitutive relation.

\section{Flow-Only Proxy Based On Volumetric Flow}

\subsection{Why Volumetric Flow Instead Of Mass Flow}

If density changes during cooldown, then mass flow changes partly because the
fluid is denser, even if the actual volume displaced by the pump has changed
less. For that reason, a flow-only viscosity proxy should be based on
volumetric flow $Q$, not mass flow $\dot{m}$.

\subsection{Definition}

The current flow-only proxy is anchored to the 3M kinematic viscosity at
$25~^\circ\mathrm{C}$:
\begin{equation}
    \nu_{\mathrm{25}} = \nu_{\mathrm{3M}}(25^\circ\mathrm{C}).
\end{equation}

The measured flow during the cooldown is compared against a fitted
$25~^\circ\mathrm{C}$ reference flow $Q_{25}$, giving
\begin{equation}
    \nu_{Q}(T) = \nu_{\mathrm{25}} \frac{Q_{25}}{Q(T)}.
\end{equation}

Dynamic viscosity then follows from
\begin{equation}
    \mu_Q(T) = \nu_Q(T)\frac{\rho(T)}{1000}.
\end{equation}

\subsection{Interpretation}

This model encodes a simple idea: for the same hardware and roughly similar
operating condition, lower volumetric flow suggests higher viscosity.

\subsection{Scientific Status}

This is a deliberately weak proxy. It ignores the full pump curve, pressure
feedback, internal recirculation, and any changes in control behavior. It is
useful as a baseline or sanity check, but it should not be interpreted as a
rigorous hydraulic inversion.

\section{Combined Pump-Flow Law}

\subsection{Definition}

The combined pump-flow law is a multivariate empirical regression. It is
designed to rely on signals that are as close as possible to viscosity-driven
hydraulic behavior, while avoiding direct use of density and avoiding absolute
system pressures that drift with temperature.

In the current ORCA implementation, the predictors are:
\begin{align}
    x_1 &= Q = \text{volume flow}, \\
    x_2 &= \Delta p_{\mathrm{pump}} = p_{\mathrm{after}} - p_{\mathrm{before}}, \\
    x_3 &= \frac{P_{\mathrm{in}}}{Q} = \text{pump input power per unit volume flow}, \\
    x_4 &= \frac{I_{\mathrm{out}}}{Q} = \text{pump output current per unit volume flow}.
\end{align}

This choice is deliberate:
\begin{itemize}
    \item $Q$ comes from the flow meter and does not use density,
    \item $\Delta p_{\mathrm{pump}}$ uses only the pressure rise across the
          pump, so common-mode pressure drift is largely removed,
    \item $P_{\mathrm{in}}/Q$ and $I_{\mathrm{out}}/Q$ act as specific loading
          indicators rather than absolute electrical signals.
\end{itemize}

So the model intentionally excludes:
\begin{itemize}
    \item direct density input,
    \item mass flow,
    \item and absolute pressure channels such as tank pressure, suction
          pressure, or discharge pressure taken separately.
\end{itemize}

The intent is not to remove all temperature sensitivity. That is impossible in
practice because viscosity itself is temperature dependent. The intent is to
remove predictors that primarily reflect density changes or common-mode system
pressure drift, so that the remaining predictors respond more directly to the
hydraulic consequences of changing viscosity.

Each feature is standardized:
\begin{equation}
    z_i = \frac{x_i - \bar{x}_i}{s_i},
\end{equation}
where $\bar{x}_i$ is the feature mean over the cooldown hold and $s_i$ is the
corresponding standard deviation.

The regression model is then
\begin{equation}
    \ln \nu_{\mathrm{comb}} = b_0 + \sum_{i=1}^{4} b_i z_i.
\end{equation}

So the predicted kinematic viscosity is
\begin{equation}
    \nu_{\mathrm{comb}} = \exp\!\left(b_0 + \sum_{i=1}^{4} b_i z_i\right),
\end{equation}
and the corresponding dynamic viscosity is
\begin{equation}
    \mu_{\mathrm{comb}} = \nu_{\mathrm{comb}}\frac{\rho}{1000}.
\end{equation}

This means the law is \emph{viscosity-focused} in its predictors, but not
strictly density-free in its final dynamic-viscosity output. Density still
appears at the last step when converting kinematic viscosity to dynamic
viscosity.

\subsection{Training Target}

The combined law is not fitted against a direct viscometer measurement. It is
trained against the local 3M-derived reference:
\begin{equation}
    \nu_{\mathrm{ref}}(T).
\end{equation}

So the fitted coefficients should be interpreted as the best local mapping from
the chosen hydraulic and electrical features to the reference viscosity over
the stable cooldown hold of the current dataset.

\subsection{Interpretation}

This model tries to infer viscosity from the way the pump and the measured
volume flow respond together during cooldown. The underlying idea is that
viscosity affects flow resistance, pump differential pressure, and electrical
loading. By using volume flow and pump differential signals instead of density
and absolute pressures, the model is made more viscosity-focused.

\subsection{Scientific Status}

This is not a first-principles hydraulic law. It is a data-driven surrogate
trained to reproduce the local 3M-derived reference viscosity from combined
pump and flow-meter signals. Its coefficients are therefore run-dependent and
operating-point dependent. It is scientifically reasonable as a multivariate
empirical model, but it should not be described as a universal pump
correlation.

\subsection{Why This Is Better Than The Earlier Pump-Only Law}

The earlier pump-only law used a wider set of pump telemetry, including
absolute pressure channels. That made it easier for the regression to absorb
temperature-driven system-pressure drift that is not directly tied to
viscosity.

The combined pump-flow law is more constrained:
\begin{itemize}
    \item it uses flow-meter volume flow instead of mass flow,
    \item it uses pump differential pressure instead of absolute pressure,
    \item and it normalizes electrical loading by flow where possible.
\end{itemize}

That does not make it a fundamental viscosity law, but it does make it a
cleaner surrogate for viscosity-driven hydraulic behavior.

\section{Combined-Law Mean Trend}

The raw combined-law prediction can be noisy because the telemetry signals are
noisy. For plotting only, a centered rolling mean is applied after sorting the
data by temperature:
\begin{equation}
    \nu_{\mathrm{comb,trend}} = \mathrm{rolling\ mean}\!\left(
    \nu_{\mathrm{comb}}
    \right).
\end{equation}

This smoothed curve is a visualization aid. It is not a separate fitted model,
and it should not replace the raw combined-law output in comparison metrics unless
that change is made explicitly and consistently.

\section{Recommended Language}

Table~\ref{tab:summary} summarizes the intended interpretation of each model.

\begin{table}[h]
    \centering
    \begin{tabular}{>{\raggedright\arraybackslash}p{3.1cm} >{\raggedright\arraybackslash}p{3.1cm} >{\raggedright\arraybackslash}p{4.1cm} >{\raggedright\arraybackslash}p{4.3cm}}
        \toprule
        Model & Primary input & Best interpretation & Main limitation \\
        \midrule
        Local smooth law & 3M viscosity table + temperature & Reference curve for the measured temperature window & Should not be extrapolated too far outside the fitted range \\
        Density-power fit & Density & Local empirical surrogate for viscosity in the measured liquid-phase range & Not a universal density-viscosity law \\
        Flow-only proxy & Volumetric flow & Very simple hydraulic baseline & Ignores most pump and pressure physics \\
        Combined pump-flow law & Volume flow + pump differential and flow-normalized electrical loading & Multivariate empirical surrogate with reduced sensitivity to density and common pressure drift & Dataset-specific and not first-principles \\
        Combined-law mean trend & Smoothed combined law & Visual guide to the average combined-law trend & Not a distinct physical model \\
        \bottomrule
    \end{tabular}
    \caption{Recommended interpretation of the current viscosity models.}
    \label{tab:summary}
\end{table}

For reports and presentations, the following phrasing is recommended:
\begin{itemize}
    \item \textbf{Local smooth law}: ``local 3M-derived reference law''
    \item \textbf{Density-power fit}: ``local empirical density-viscosity surrogate''
    \item \textbf{Flow-only proxy}: ``volume-flow-based viscosity proxy''
    \item \textbf{Combined pump-flow law}: ``combined pump-flow viscosity surrogate''
\end{itemize}

\section{Practical Recommendation}

For the current ORCA workflow, the most conservative hierarchy is:
\begin{enumerate}
    \item use the local 3M-derived law as the primary reference,
    \item use the density-power fit as the strongest local surrogate,
    \item use the combined pump-flow law as a telemetry-based secondary surrogate,
    \item use the flow-only proxy only as a simple baseline.
\end{enumerate}

This hierarchy matches both the physics and the current data quality: density
is usually the strongest single predictor in the measured cooldown hold, while
the combined pump-flow and flow-only models remain useful mainly as supporting
comparisons.

\end{document}
